\noindent

The Light Dark Matter eXperiment (LDMX) searches for sub-GeV dark matter
through missing-momentum measurements. Multiple electrons in one readout
window can produce overlapping electromagnetic-calorimeter showers and
obscure hit origins. This thesis studies supervised hit-level separation in
simulated two- and three-electron pile-up events.

Events are represented as unordered, variable-size sets of calorimeter hits
and trigger-pad-track observations. Transformer and GravNet architectures
assign each hit to its dominant origin electron; a multitask Transformer also
predicts electron count, per-hit contributor sets, and deposited-energy
fractions.

On matched test samples, the Transformer achieved pooled hit accuracies of
\(83.31\,\%\) for two-electron events and \(73.11\,\%\) for three-electron
events, compared with \(83.09\,\%\) and \(72.63\,\%\), respectively, for
GravNet. The architectures therefore performed similarly, with a small
Transformer advantage. Three-electron events were harder because of closer
shower overlap, while higher energy-weighted accuracies showed that
misassigned hits tended to carry slightly less energy. Trigger-pad information had
negligible measured influence for fixed multiplicity but improved
variable-multiplicity electron counting. The final model achieved
electron-count, exact contributor-set, and dominant-origin accuracies of
\(99.98\,\%\), \(72.41\,\%\), and \(77.04\,\%\), respectively.

Set-based neural architectures can thus recover substantial per-electron
shower structure from simulated LDMX pile-up events, although mixed and
strongly overlapping deposits remain the principal challenge. The reusable
framework enables further studies with heterogeneous LDMX detector inputs.




























%\textbf{Baseline: Supervised hit-to-particle clustering}

%\textbf{"MLPF-lite": Multi-task hit-level shower reconstruction}

%\noindent 
%Imagine matter, like the bread you eat every morning, but now imagine matter but dark, but not like dark bread that you know of but instead invisible dark matter bread that you cannot see, but do not dare think that the matter does not matter even though the matter cannot be seen. Dark matter matters because it is matter. This thesis matters because it attempts to make a ML reconstruction method that matter in a dark matter experiment. Yes!

%\lipsum[1]